% WITHDRAWN 2026-11-19 08:16 UTC. DO NOT INCLUDE IN THE REPORT.
% The proposed degrees 4,5,6 computation incorrectly interchanged
% (Sym^2 V)^vee with Sym^2(V^vee) in characteristic two.
% In fact Frobenius splitting gives Sym^2 V = V directsum L2^20,
% and the nonsplit Tango class implies H0(V^vee)=0.
% The claimed threshold 7 is therefore unsupported and the r6=1 claim false.
% This file is retained solely as a transparent record of the withdrawn route.

\subsection{The exact injectivity threshold for the explicit threefold}

\begin{lemma}[Negative first cohomology in the two-step example]
\label{lem:OC-threefold-negative-degrees}
\status{Candidate proof; independent audit pending.}
Keep $S,L_2,R,L,X,A$ from
Proposition~\ref{prop:OC-Tango-threefold} and
Lemma~\ref{lem:OC-threefold-multiroot}.  Then
\[
 h^1(R,L^{-4})=69,\qquad h^1(R,L^{-5})=24,\qquad
 h^1(R,L^{-6})=1,\qquad H^1(R,L^{-r})=0\quad(r\geq7).
\]
Consequently
\[
 h^1(X,A^{-4})=376{,}047,\qquad h^1(X,A^{-5})=867,
 \qquad h^1(X,A^{-6})=1,
 \qquad H^1(X,A^{-d})=0\quad(d\geq7),
\]
while $h^1(X,A^{-d})\geq843$ for every $1\leq d\leq5$.
\end{lemma}
\begin{proof}
For the second normalized degree-five cover, put
\[
 m_j(r)=\left\lceil\frac{2(r+j)}5\right\rceil,
 \qquad e_j(r)=4j-r-10,\qquad0\leq j\leq4.
\]
The negative eigensheaves are
\[
 \psi_{2*}L^{-r}
  =\bigoplus_{j=0}^4\OO_{P_2}(-m_j(r))\otimes L_2^{4j-r}.
\]
For $r\geq1$ all relative degrees are strictly negative, so $g_*L^{-r}=0$.  Relative duality and Leray identify $H^1(R,L^{-r})$ with the sum, over indices having $m_j(r)\geq2$, of
\begin{equation}\label{eq:OC-negative-R1}
 H^0\!\left(S,\operatorname{Sym}^{m_j(r)-2}(V_2)^\vee
                                  \otimes L_2^{e_j(r)}\right).
\end{equation}
Here $V_2$ fits into $0\to\OO_S\to V_2\to L_2^{10}\to0$, so the indicated symmetric power has a filtration with quotients $L_2^{-10a}$.
If $r\geq7$, every $e_j(r)\leq6-r<0$.  Every filtered line quotient is antiample, proving the asserted vanishing for all such $r$, without a finite-range inference.

For $r=4,5,6$, every index except $j=4$ again has only negative filtered quotients.  At $j=4$ one has $m_j(r)=4$ and the contribution is
\[
        H^0(S,\operatorname{Sym}^2(V_2)^\vee\otimes L_2^{6-r}).
\]
In characteristic two the square map gives a subbundle
\[
 F_S^*(V_2^\vee)\hookrightarrow\operatorname{Sym}^2(V_2^\vee).
\]
The Tango extension is Frobenius-split, hence the left side is
$\OO_S\oplus L_2^{-20}$.  This embeds all sections of $L_2^{6-r}$ into the preceding group.  Conversely its filtration has quotients
$L_2^{6-r}$, $L_2^{-4-r}$, and $L_2^{-14-r}$, the latter two antiample.  Its number of sections is therefore exactly $h^0(S,L_2^{6-r})$.

The first-cover formula gives
\[
 f_*L_2=N_0\oplus N_0^{21},\qquad
 f_*L_2^2=N_0^2\oplus N_0^{22}\oplus N_0^{42}.
\]
For $0\leq b\leq42$, $h^0(C,N_0^b)=b+1$.  Indeed a function with its only pole at infinity is in
$k[x,y]/(y^2+y-x^{843})$; writing it uniquely as $a(x)+y b(x)$, the two summands have even and odd pole orders, and every nonzero second summand has pole order at least $843$.  Below this bound the sections are exactly $1,x,\ldots,x^b$.  Thus
\[
       h^0(L_2)=2+22=24,\qquad h^0(L_2^2)=3+23+43=69.
\]
Also $h^0(\OO_S)=1$, proving the three asserted $H^1$ dimensions.

On the multiroot cover,
\[
 H^1(X,A^{-d})=
  \bigoplus_{0\leq a_i<445}H^1(R,L^{-(d+\sum a_i)}).
\]
For $d\geq4$, only sums $\sum a_i\leq6-d$ contribute.  There are $843$ tuples of sum one and $\binom{844}{2}=355{,}746$ tuples of sum two.  This gives
\[
 69+843\cdot24+355{,}746=376{,}047,\quad
 24+843=867,\quad1,
\]
for $d=4,5,6$ respectively, and zero thereafter.
For $d=1$ use the $843$ tuples of sum one and $H^1(R,L^{-2})\ne0$.  For $d=2$ use the $843$ tuples with one entry equal to two and the exponent-four group; for $d=3$ use the tuples of sum one and the exponent-four group.  These observations prove all the lower bounds through degree five.
\end{proof}

\begin{lemma}[A cotangent vanishing bound with no positivity estimate left implicit]
\label{lem:OC-threefold-differential-degrees}
\status{Candidate proof; independent audit pending.}
For the same example,
\[
 H^0(R,\Omega_R^1\otimes L^{-r})=0\quad(r\geq5),
 \qquad
 H^0(X,\Omega_X^1\otimes A^{-d})=0\quad(d\geq5).
\]
For every smooth $Y\in|A^d|$, $d\geq5$, restriction
$H^0(X,\Omega_X^1)\to H^0(Y,\Omega_Y^1)$ is injective.
\end{lemma}
\begin{proof}
The first Tango step supplies a nowhere-zero normalized differential, equivalently a subbundle
\[
                L_2^{20}\hookrightarrow\Omega_S^1.
\]
Let $Q$ be its locally free quotient of rank one.  On
$P_2=\mathbf P_S(V_2)$, the logarithmic cotangent bundle for the disjoint smooth divisors $E_0,G$ has a filtration with successive quotients
\begin{equation}\label{eq:OC-log-cotangent-filtration}
 \OO_{P_2}(2),\qquad \pi^*Q,\qquad
 \Omega^1_{P_2/S}(E_0)=\OO_{P_2}(-1)\otimes L_2^{10}.
\end{equation}
We check the filtration directly.  On an affine fibre chart $G$ has equation $S^2-f_i=0$, where $df_i$ generates the indicated base subbundle.  The only new horizontal logarithmic generator is $df_i/(S^2-f_i)$.  Thus the horizontal kernel is an extension of $\pi^*Q$ by
\[
 \pi^*L_2^{20}(G)=\OO_{P_2}(2),
\]
since $\OO(G)=\OO(2)\otimes L_2^{-20}$.  The relative quotient gains precisely the simple pole along the infinity section $E_0$, and no pole along $G$, because $d_{P_2/S}(S^2-f_i)=0$.  Near infinity $G$ is absent; replacing $S$ by $z^{-1}$ gives the same horizontal subbundle and the relative generator $dz/z$.  These local descriptions glue and give \eqref{eq:OC-log-cotangent-filtration}.

The normalized cover is tame and logarithmically \'etale along these branches: its local equations are $s=u^5$ and, along the other branch, a fifth root of a smooth parameter.  Thus
\[
           \Omega_R^1\hookrightarrow
           \psi_2^*\Omega_{P_2}^1(\log(E_0+G)).
\]
Tensor with $L^{-r}$, push down, and use the normalized eigensheaves from the preceding lemma.  The three corresponding graded terms for index $j$ are
\[
 \OO_{P_2}(2-m_j(r))\otimes L_2^{4j-r},\qquad
 \OO_{P_2}(-m_j(r))\otimes\pi^*Q\otimes L_2^{4j-r},
\]
\[
 \OO_{P_2}(-1-m_j(r))\otimes L_2^{4j-r+10}.
\]
The last two have no sections because their relative degrees are negative.  If $r\geq6$, then every $m_j(r)\geq3$, and the first also has negative relative degree.  At $r=5$, only $j=0$ has $m_j(r)=2$; its first term is $L_2^{-5}$, with no sections.  This proves the first vanishing for every $r\geq5$.

For the $843$-root cover with branch divisor $D=\sum D_i$, the same local logarithmic identification gives
\[
 \Omega_X^1\hookrightarrow\rho^*\Omega_R^1(\log D).
\]
Its eigensheaves after twisting by $A^{-d}$ are bounded by
$\Omega_R^1(\log D)\otimes L^{-(d+\sum a_i)}$.
For every positive $r$, the residue exact sequence shows
\[
 H^0(R,\Omega_R^1(\log D)\otimes L^{-r})
     =H^0(R,\Omega_R^1\otimes L^{-r}):
\]
each residue target $H^0(D_i,L^{-r}|_{D_i})$ is zero, as the restricted bundle is antiample.  All the indices here satisfy $r=d+\sum a_i\geq5$, proving the second vanishing.

Finally the restriction sequence and conormal sequence show that the kernel on global one-forms is $H^0(X,\Omega_X^1(-Y))$: the additional conormal contribution has no global sections since
$H^0(Y,\OO_Y(-Y))=0$.  The second vanishing therefore proves the last assertion.
\end{proof}

\begin{lemma}[The integral first crystalline lattice comes from the base curve]
\label{lem:OC-integral-first-lattice}
\status{Candidate proof; independent audit pending.}
Let $h:X\to C$ be the composite map in the explicit construction.  Then $h_*\OO_X=\OO_C$, its geometric generic fibre is integral, and pullback is an integral isomorphism
\[
                 H^1_{\cris}(C/W)\xrightarrow{\sim}H^1_{\cris}(X/W).
\]
For every smooth section $Y\in|A^d|$ with $d\geq5$, the induced morphism $h_Y:Y\to C$ satisfies $(h_Y)_*\OO_Y=\OO_C$, and pullback gives integral isomorphisms
\[
       H^1_{\cris}(C/W)\xrightarrow{\sim}H^1_{\cris}(Y/W),
       \qquad H^1_{\cris}(X/W)\xrightarrow{\sim}H^1_{\cris}(Y/W).
\]
\end{lemma}
\begin{proof}
First, each Tango step is flat with geometrically integral fibres, as follows from its finite locally free normalized algebra and its explicit cuspidal fibre.  Their composite $R\to C$ therefore has geometrically integral fibres.  To see integrality directly, base-change to a geometric point of $C$: the base cusp is integral, the second morphism is flat, and its geometric generic fibre is integral.  Any affine chart injects into its generic fibre algebra by flatness over the base domain, so it is reduced; all nonempty open charts meet the same integral generic fibre, giving irreducibility.  Hence $k(R)/k(C)$ is a regular field extension.

Consider the Stein factorization $X\to C'\to C$.  Because $X$ is normal, $C'$ is a smooth projective curve.  Its function field is the algebraic closure of $k(C)$ in $k(X)$.  Regularity of $k(R)/k(C)$ makes it linearly disjoint from $k(C')$; consequently
\[
                  [k(C'):k(C)]\mid[k(X):k(R)]=445^{843}.
\]
The finite cover $C'\to C$ is therefore separable.  Stein factorization gives geometrically connected fibres for $X\to C'$, so proper base change and Leray inject $H^1_{\et}(C',\mathbf Q_3)$ into $H^1_{\et}(X,\mathbf Q_3)$.  Its dimension is at most $842=2g(C)$.  Riemann--Hurwitz, including its effective different in the possibly wild case, gives
\[
                  g(C')-1\geq\deg(C'/C)\,(421-1).
\]
Thus $\deg(C'/C)=1$.  This proves $h_*\OO_X=\OO_C$.
Also $k(X)/k(C)$ is separably generated: $k(R)/k(C)$ is regular and $k(X)/k(R)$ is separable.  Since it has no new algebraic constants, it is regular, proving that the geometric generic fibre is integral.  Exact primary locators are
\cite[Theorem~37.53.4, Tag~03H0; Lemma~53.12.2, Tag~0C1D, in Section~53.12, Tag~0C1B]{StacksProject}.

We now identify the integral lattice, not just its rational span.  Both Tango steps have sections.  Their pullbacks on reduced Picard varieties are split monomorphisms; the first Betti numbers are equal, so these monomorphisms are isomorphisms.  Thus
\[
                   \Pic^0_{R,\mathrm{red}}\simeq\Pic^0_C.
\]
For the finite locally free cover $\rho$, the norm satisfies
\[
            \operatorname{Nm}_\rho\circ\rho^*=[445^{843}].
\]
The norm and this identity commute with arbitrary base change, by
\cite[Lemmas~31.18.2 and~31.18.6, Tags~0BCY and~0BD2]{StacksProject}; hence they give morphisms of Picard group schemes.  The reduced Picard varieties of $R$ and $X$ have the same dimension $421$.  Their pullback is therefore an isogeny whose kernel is contained in the torsion killed by the odd integer $445^{843}$.  It induces an isomorphism of their $2$-divisible groups.  Illusie's integral functorial identification of first crystalline cohomology with the reduced Picard realization
\cite[II, Proposition~3.11 and Remark~3.11.2]{Illusie1979DRW}
proves the first integral isomorphism in the statement.  This is where the odd covering degree, rather than rational Betti equality alone, is used.

Let now $Y$ be any smooth section of degree $d\geq5$.  It dominates $C$: an ample divisor cannot be disjoint from a positive-dimensional fibre.  Over an algebraic closure of $k(C)$, it is an effective ample Cartier divisor in the integral projective surface $X_{\overline\eta}$.  This surface is Cohen--Macaulay: it is a geometric fibre of a flat morphism from a smooth variety to a smooth curve, hence locally a Cartier complete intersection.  An ample effective divisor in such a surface is connected.  Here is the precise cohomological argument: for large $n$, Serre duality and Serre vanishing give
\[
               H^1(X_{\overline\eta},\OO(-nY))=0.
\]
Since $H^0(X_{\overline\eta},\OO)=\overline{k(C)}$, the divisor sequence gives the same constants on $nY$.  If $Y_{\overline\eta}$ were disconnected, every $nY$ would have a nontrivial idempotent global function, a contradiction.  Thus the geometric generic fibre of $Y\to C$ is connected.

It is also generically separable.  A nonzero global differential on $C$ pulls back nontrivially to $X$, since $k(X)/k(C)$ is regular, and then nontrivially to $Y$ by Lemma~\ref{lem:OC-threefold-differential-degrees}.  The map from the one-dimensional $k(C)$-space $\Omega^1_{k(C)/k}$ to $\Omega^1_{k(Y)/k}$ is consequently nonzero, hence injective; the separating-transcendence-basis criterion gives separable generation of $k(Y)/k(C)$.  In its Stein factorization, geometric generic connectedness makes the finite curve map purely inseparable, while separable generation makes it separable.  It is therefore the identity, and $(h_Y)_*\OO_Y=\OO_C$.

This equality is enough to make $\Pic_C\to\Pic_Y$ a monomorphism of fppf Picard sheaves.  Indeed it remains true after every base change over $k$, by flat base change.  If the pullback of a line bundle from $C$ is trivial, or comes from the parameter base, projection formula and pushforward recover the same assertion on $C$.  The same reasoning after a local fppf trivialization handles the sheaf quotient defining the Picard functor.  Restriction to the reduced identity components is thus a closed immersion of abelian varieties.

Finally $b_1(Y)=b_1(X)=842$.  One can use the usual \'etale weak Lefschetz statement, or obtain it directly here: $X\setminus Y$ is smooth affine of dimension three, so affine vanishing and Poincar\'e duality give $H_c^1=H_c^2=0$; localization then identifies $H^1(X,\mathbf Q_3)$ and $H^1(Y,\mathbf Q_3)$, with the same SGA~4 locators as in Lemma~\ref{lem:OC-threefold-multiroot}.  The closed immersion $\Pic_C^0\to\Pic^0_{Y,\mathrm{red}}$ is between abelian varieties of dimension $421$, hence is an isomorphism.  Illusie's integral identification yields the first isomorphism for $Y$.  Factoring that map through $X$ yields the final isomorphism of first crystalline lattices.
\end{proof}

\begin{theorem}[Exact degree range for integral injectivity in the example]
\label{thm:OC-exact-integral-injectivity-degrees}
\status{Candidate proof; independent audit pending.}
For the fixed embedded smooth threefold $(X,A)$ of
Lemma~\ref{lem:OC-threefold-multiroot} and every smooth degree-$d$ hypersurface section $Y_d$, the map
\[
            H^2_{\cris}(X/W)\longrightarrow H^2_{\cris}(Y_d/W)
\]
is injective if and only if $d\geq7$.  Its kernel has a nonzero subgroup killed by two for every $1\leq d\leq6$.
More precisely, in degrees five and six the dimensions of those subgroups are
\[
 \dim_k\ker\bigl(H^2_{\cris}(X/W)[2]\to H^2_{\cris}(Y_d/W)[2]\bigr)
   =\begin{cases}867,&d=5,\\1,&d=6.\end{cases}
\]
\end{theorem}
\begin{proof}
For $1\leq d\leq5$, Lemma~\ref{lem:OC-threefold-negative-degrees} gives a coherent restriction kernel of dimension at least $843$.  The conjugate-edge and UCT argument in Theorem~\ref{thm:OC-integral-injectivity-counterexample} applies unchanged, because the first crystalline rank remains $842$.  Thus restriction fails to be injective throughout these degrees.

For all $d\geq5$, we can compute the kernel more precisely.  The low-degree conjugate exact sequence starts
\[
 0\to H^1(X^{(1)},\OO_{X^{(1)}})
   \to H^1_{\dR}(X/k)
   \to H^0(X^{(1)},\Omega^1_{X^{(1)}})
   \to H^2(X^{(1)},\OO_{X^{(1)}}).
\]
The one-form restriction is injective by
Lemma~\ref{lem:OC-threefold-differential-degrees}, also after Frobenius twist.  Functoriality therefore identifies the de Rham restriction kernel with the coherent restriction kernel under the conjugate edge.  Since the smooth ample section is connected, the latter is $H^1(X,A^{-d})^{(1)}$.

Compare the two natural universal-coefficient sequences
\[
 0\to H^1_{\cris}(-/W)/2\to H^1_{\dR}(-/k)
               \to H^2_{\cris}(-/W)[2]\to0
\]
for $X$ and $Y_d$.  Their left map is an isomorphism by
Lemma~\ref{lem:OC-integral-first-lattice}.  The snake lemma consequently identifies the right-hand kernel with the de Rham kernel.  Hence, for every $d\geq5$,
\[
 \dim_k\ker\bigl(H^2_{\cris}(X/W)[2]\to H^2_{\cris}(Y_d/W)[2]\bigr)
                       =h^1(X,A^{-d}).
\]
The preceding negative-cohomology lemma gives $867$ in degree five, $1$ in degree six, and zero for every $d\geq7$.

For $d\geq7$, zero kernel on the subgroup killed by two implies zero full kernel: rational crystalline weak Lefschetz makes the full kernel a finite torsion module, and every nonzero finite $2$-primary torsion module has a nonzero element killed by two.  The exact primary rational comparison used here is
\cite[Appendix~9, Theorem~9.1 and Corollary~9.3, pp.~63--64, arXiv version~3]{NakkajimaYobuko2022}; its proof uses the compact-support rigid sequence, duality, and the length of the affine Monsky--Washnitzer complex.  This proves both directions of the claimed threshold.
\end{proof}

\begin{remark}[Scope of the degree calculation]
\label{rem:OC-exact-degree-scope}
This computes the full degree range for injectivity A.1(c) on this fixed example, and the full coherent first-cohomology restriction range as well.  It does not identify the cokernel torsion or prove an isomorphism of the full degree-two torsion subgroups.  The separate exact-degree questions for A.1(a) and A.1(b), including the earlier product examples, are not replaced by an assertion about subgroups killed by two.
\end{remark}
